% Sumário executivo. Deliberadamente NÃO numerado (\section*), para não
% deslocar a numeração das demais seções, que é referida em comentários do
% principal.tex e em conversas do projeto. Entra no índice à mão.
%
% Linguagem (revisão de 01/08/2026, a pedido de João): este é o documento
% que o gestor lê primeiro e, muitas vezes, o único que ele lê. O texto foi
% reescrito para falar de operação, capacidade e resultado, e não de método:
% o nome do estimador, o erro-padrão e o valor de p ficam nas seções, que é
% onde quem audita vai procurá-los. Nada foi arredondado nem suavizado no
% caminho; o que mudou foi o vocabulário, não o número.
\section*{Sumário executivo}
\addcontentsline{toc}{section}{Sumário executivo}
\label{sec:resumo}

\noindent Este documento compara o desempenho dos cinco arranjos
institucionais que administram hospitais do SUS no estado de São Paulo,
sobre um painel de \textbf{289 hospitais acompanhados ano a ano de 2015 a
2025} (3.179 observações, sem buracos), destilado de 830 estabelecimentos
da base bruta do SIH/SUS. As categorias são Filantrópico, Público
Municipal, gestão por Organização Social de Saúde (OSS), administração
Direta e Privado.

\medskip
\noindent\textbf{A resposta curta é que a gestão por Organização Social
lidera, e lidera com folga, nas duas dimensões em que a evidência é
firme.} As OSS entregam mais atendimento com a mesma estrutura e
desperdiçam menos recurso no caminho. Os dois resultados se mantêm em
todos os testes a que foram submetidos, inclusive nos que retiram da
amostra os hospitais que trocaram de gestão no período, e por isso não
são efeito de um caso isolado nem artefato de um método.

\medskip
\noindent\textbf{A pergunta, porém, se decompõe em três}, e a liderança
precisa ser lida com precisão: quem opera com mais intensidade, quem
transforma recurso em atendimento com menos desperdício e quem
\emph{provoca} melhor resultado clínico ao mudar de gestão são perguntas
diferentes. As duas primeiras estão respondidas, e a favor das OSS. A
terceira repousa sobre cinco hospitais convertidos e admite no máximo
evidência sugestiva. O que segue está ordenado pela firmeza da evidência,
não pela ordem dos métodos.

\subsection*{O que se sustenta}

\textbf{As OSS operam num ritmo que nenhuma outra rede alcança.}
Comparados hospitais de mesma complexidade, mesmo porte e mesmo ano, as
OSS \textbf{atendem 43,7\% mais pacientes}, com os leitos de internação
\textbf{22,7\% mais ocupados} e a UTI \textbf{19,3\% mais usada}, entre os
hospitais que a possuem (Seção~\ref{sec:resultados}). A segunda fase da
análise, que compara apenas hospitais de capacidade instalada
equivalente, mede a intensidade de UTI \textbf{30,4\% acima da
administração Direta} (Seção~\ref{sec:fase2}). E a vantagem não depende de
dinheiro nem de escolha de método: em produção física por leito, que não
usa valor nenhum, são \textbf{55 saídas por leito nas OSS contra 36 na
Direta}, ou seja, cada leito da rede OSS atende meia vez mais gente por
ano.

\textbf{As OSS desperdiçam menos: são a rede que mais se aproxima da
melhor prática do estado.} Medida a distância de cada hospital até a
fronteira das melhores práticas observadas, a nota mediana do período é de
\textbf{0,733 nas OSS}, contra 0,634 na administração Direta, 0,632 no
Público Municipal e 0,595 no Filantrópico. São \textbf{três patamares, e
não quatro}: Direta e Município empatam no meio. A vantagem das OSS
equivale a cerca de \textbf{25\% menos distância até a fronteira}, e
sobrevive à retirada dos hospitais convertidos
(Seção~\ref{sec:fase2}). Traduzido para a decisão de quem administra: com
a mesma estrutura física e o mesmo quadro, a rede OSS entrega mais
atendimento e deixa menos capacidade parada.

\textbf{Os dois traços acima são características das redes}, medidas entre
hospitais distintos, e valem como retrato de quem opera melhor hoje. Elas
não são, por si, prova de que converter um hospital produz esse resultado:
essa é a terceira pergunta, e é a que o dado disponível não fecha.

\subsection*{O que é sugestivo, e o que não se sustenta}

\textbf{Sugestivo:} o paciente fica cerca de \textbf{10\% menos tempo
internado} sob gestão por OSS, medido dentro do mesmo hospital, antes e
depois da conversão. O resultado resiste aos testes de robustez que o
número de hospitais convertidos permite fazer, e não é contrariado por
nenhum deles; os três testes que não o confirmam são justamente os que
perdem força quando há apenas cinco hospitais tratados. É evidência
sugestiva forte, e não fato demonstrado
(Seções~\ref{sec:efeitogestao} e~\ref{sec:inferenciafinal}).

\textbf{Não se sustenta:} a queda de mortalidade depois da conversão. Ela
aparece nos dados brutos, some quando cada hospital convertido é comparado
a um hospital gêmeo construído a partir dos não convertidos, e desaparece
por inteiro quando os cinco conversores saem da amostra. \textbf{Nada
neste documento autoriza afirmar que a gestão por OSS reduz a
mortalidade} (Seções~\ref{sec:fase2} e~\ref{sec:conclusao}).

\textbf{Sem diferença em método nenhum:} o faturamento por saída. Aqui
cabe uma ressalva de leitura importante para quem for usar o número em
negociação de contrato: as OSS são custeadas por contrato de gestão, e não
por reembolso de internação, de modo que o valor por saída pode estar
refletindo o regime de pagamento e não o consumo de recursos.

\subsection*{Duas descobertas sobre os próprios números}

\textbf{Hospital com ocupação acima de 100\% não é erro de cadastro: é
hospital superlotado.} Isso acontece em 15,5\% das observações, e a
explicação é simples. O cadastro nacional só registra o leito oficialmente
habilitado; já o sistema que paga a internação paga também a diária do
leito extra que o hospital montou para dar conta da demanda. Quando o
hospital atende mais gente do que os leitos oficiais comportariam, a conta
passa de 100\%. O que está ali é informação verdadeira sobre sobrecarga, e
não sujeira de dado.

Isso obriga a uma correção de leitura. O número que o documento chama de
``ocupação de internação'' estava respondendo a \textbf{duas perguntas ao
mesmo tempo}: o quanto o hospital opera cheio, e com que frequência ele
transborda. Separadas as duas, entre hospitais de estrutura equivalente a
média das OSS é de \textbf{93,0\% contra 69,4\% da administração Direta},
uma distância de 23,5 pontos --- e dessa distância \textbf{21,5 pontos
(91\%) são leito de fato ocupado, e só 2,0 pontos (9\%) são transbordo}.
Em uma frase: a mistura existia, mas explicava menos de um décimo da
diferença, e \textbf{o ritmo mais intenso das OSS continua inteiro depois
de separá-la} (Seção~\ref{sec:segmentacao}).

\textbf{Só uma parte da rede permite comparar modelo de gestão com modelo
de gestão.} Comparação honesta exige hospitais de estrutura parecida, e por
isso os 289 foram separados em três grupos: os que não têm UTI do SUS
(67 hospitais), os que têm UTI e complexidade mais baixa (97) e os que têm
UTI e alta complexidade (125). Nos dois primeiros grupos praticamente não
existe hospital de administração Direta nem de OSS: são \textbf{1 e 6
hospitais da Direta contra 2 e 7 das OSS}. Comparar redes com esse punhado
de casos é comparar hospitais avulsos e chamar o resultado de modelo de
gestão. Sobra o terceiro grupo, com \textbf{28 hospitais de OSS contra 20
da Direta}: é o único em que a comparação se sustenta, e é exatamente ali
que a vantagem das OSS foi medida. A análise que trata a rede inteira de
uma vez não deixava essa limitação aparecer (Seção~\ref{sec:segmentacao}).

\subsection*{As quatro limitações que condicionam tudo acima}
\begin{enumerate}[nosep,leftmargin=*]
\item \textbf{A resposta causal repousa em cinco hospitais convertidos},
  três deles com um único ano de dados depois da mudança. Nenhum é caso
  limpo: uma conversão respondeu a uma crise de gestão, outra veio junto
  com prédio novo, dois hospitais já tinham serviços terceirizados antes,
  e os dois únicos com histórico longo foram para a \textbf{mesma
  organização social} (Subseção~\ref{sub:oqueforam}). É por isso que a
  liderança das OSS está afirmada como retrato da rede, e não como
  promessa de que converter produz o mesmo resultado em qualquer hospital.
\item \textbf{A eficiência é medida por um método só.} Dois foram
  calculados; o segundo falhou, devolvendo praticamente a mesma nota para
  todas as OSS ao longo de onze anos, o que é o modelo travando e não um
  retrato da rede. Como as OSS são justamente o grupo que a pesquisa
  compara, ele foi retirado das medidas e das listas
  (Apêndice~\ref{ap:sfa}).
\item \textbf{O grupo Privado tem 3 hospitais} e nunca é interpretado:
  entra nos modelos apenas para não contaminar as demais comparações.
\item \textbf{Não há ajuste pela gravidade clínica do paciente.} A base
  não informa o quanto cada caso era grave, e a complexidade foi medida
  pela estrutura física e tecnológica do hospital, que é um substituto
  imperfeito. Na prática: um hospital que recebe casos mais graves pode
  aparecer com resultado pior sem que a gestão tenha qualquer
  responsabilidade sobre isso.
\end{enumerate}

\subsection*{Uso indicado}

\noindent O que a pesquisa entrega com segurança é o \textbf{retrato do
regime operacional e da eficiência produtiva das cinco redes}, e nele as
\textbf{OSS estão à frente por margens medidas, reproduzidas em todos os
métodos e resistentes a todos os testes aplicados}. O que ela
expressamente não entrega é um veredito sobre qual modelo salva mais
vidas: essa pergunta exigiria mais hospitais convertidos e mais anos de
acompanhamento do que os dados de hoje oferecem. Os resultados foram
construídos para serem lidos em conjunto e com o contexto metodológico
completo; isolado desse contexto, qualquer número aqui fica vulnerável ao
uso indevido.
